% Generated by gallery.py from gallery.yaml and the boards repo.
% Edit those and run sync.py; never edit this file.

% ======================================================================
\section{The boards}
\hslead{Twenty-two boards have reached a routed layout. They run from a linear
regulator with five parts to a three-phase motor driver.}

\begin{hsblock}
\hslabel{Table 1 / the twenty-two boards and where each one stopped}\par\vspace{8pt}
\footnotesize
\begin{tabularx}{\linewidth}{@{}l Y r r >{\raggedright\arraybackslash}p{3.3cm}@{}}
\hstoprule
\hshead{Board} & \hshead{What it is} & \hshead{Layers} & \hshead{Size (mm)} & \hshead{Where it stopped}\\
\hstoprule
pd-trigger & USB-C PD bench trigger & 2 & 48 $\times$ 30 & ordered and made\\\hline
lumina-carrier & lighting carrier, PoE-powered & 4 & 100 $\times$ 80 & ordered and made\\\hline
rp2040-mini & RP2040 dev board on the Pico footprint, USB-C & 2 & 21 $\times$ 51 & package ready\\\hline
esp32c3-node & WiFi/BLE sensor node, LiPo & 2 & 26 $\times$ 38 & package ready\\\hline
lipo-boost & LiPo charger and 5\,V boost & 2 & 53.9 $\times$ 38.8 & verification, findings open\\\hline
nfc-card & NFC business card & 2 & 88.9 $\times$ 50.8 & package ready\\\hline
stereo-class-d-amp & stereo class-D amplifier & 2 & 53.6 $\times$ 36.8 & verification, findings open\\\hline
bldc-motor-driver & three-phase motor driver & 4 & 78 $\times$ 78 & placement, findings open\\\hline
pd-trigger-lite & small USB-C PD trigger & 2 & 25 $\times$ 15 & package ready\\\hline
pd-trigger-lite-dip & the same, set by DIP switch & 2 & 25 $\times$ 21 & package ready\\\hline
g0-sense & USB-C sensor node & 2 & 35.8 $\times$ 28.3 & package ready\\\hline
bb-adc & single-channel ADC & 2 & 54.8 $\times$ 34.9 & package ready\\\hline
bb-mcu & bare microcontroller board & 2 & 34.8 $\times$ 22.3 & package ready\\\hline
bb-ldo & bare linear regulator & 2 & 34.7 $\times$ 34.7 & package ready\\\hline
bb-amp & sensor amplifier chain & 2 & 48 $\times$ 28.3 & package ready\\\hline
bb-buck & bare buck converter & 2 & 35 $\times$ 25 & package ready\\\hline
rf-term-150w & 150\,W RF dummy load head & 2 & 26 $\times$ 20 & package ready\\\hline
rf-de-20m & 20\,MHz Class E GaN stage & 4 & 120 $\times$ 80 & routing, findings open\\\hline
sbuck-5v3a & 5\,V / 3\,A synchronous buck & 4 & 50 $\times$ 40 & package ready\\\hline
usb-buck & USB device dev board & 4 & 50 $\times$ 40 & package ready\\\hline
lumina-par & RGBW PAR daughter board & 4 & 100 $\times$ 80 & manufacturability, findings open\\\hline
stm32-blinky & minimal STM32 dev board & 2 & 50 $\times$ 40 & package ready\\
\hstoprule
\end{tabularx}
\end{hsblock}

``Where it stopped'' is the last stage the board reached, and it does not mean
the board is fit to release. Two boards were ordered and made, and their files
record the JLCPCB order numbers. ``Package ready'' means every recorded check
passed and the manufacturing package is built. ``Findings open'' means a check
failed and a person has yet to rule on it. Each size is the bounding box of
the board outline, and not every outline is a rectangle.

\hsprovenance{Generated by docs/showcase/gallery.py from each board's .kicad\_pcb (layers,
size, footprints), state.json and fab folder (stage) in the boards repository,
as last changed on 4 October 2026.}
\newpage

% ----------------------------------------------------------------------
\subsection{New boards}

These boards were added after the rest of this chapter was written, so each
description is taken from the board's own requirements.

\boardpair
{\board{rp2040-mini}{rp2040-mini}%
{2 layers \sep 21 $\times$ 51\,mm \sep 53 footprints}{package ready}%
{A general-purpose RP2040 microcontroller board that is a drop-in for the
Raspberry Pi Pico (Raspberry Pi Ltd, 2021): same 21 x 51\,mm outline, same
40-pin through-hole edge (as built: plain PTH, not castellated) with the
Pico's pinout, same 3-pin SWD header position class (bottom edge, opposite
USB).}}
{}

\hsprovenance{Pictures from docs/showcase/renders, made as that folder's README says. Facts
from the board files of rp2040-mini in the boards repository.}
\newpage

% ----------------------------------------------------------------------
\subsection{The motor driver}
\label{pg:bldc}

\board{bldc-motor-driver}{bldc-motor-driver}%
{4 layers \sep 78 $\times$ 78\,mm \sep 175 footprints}{placement, findings open}%
{A three-phase driver for brushless motors, running from 10 to 28\,V. Six
N-channel MOSFETs make three half-bridges, a DRV8300 drives their gates, and
an STM32G431 runs the motor control. Each low-side leg has a 3\,m$\Omega$
shunt read by an INA240 current amplifier. The brief was a general spec and
one frame of a video showing someone else's demo board, so hwde wrote down
every gap it had to fill as a numbered assumption.}

The FETs are rated 60\,V although the board runs at 28\,V at most. The input
TVS diode clamps a surge at about 45\,V, and every part on the motor supply is
chosen to survive that clamp, not just the working voltage. There is no fuse
on the board, and the notes call for a 20\,A slow fuse in the supply lead.

\newpage

\board{bldc-motor-driver-fcu}{bldc-motor-driver, top copper}%
{F.Cu with net names}{}%
{The three half-bridges sit along the bottom edge, with phase A, B and C left
to right, each with its current amplifier below it. The board passes its
checks with 390 written waivers, and the board's notes ask you to read them
before you trust it. Most cover stitching vias and the thin Kelvin and
gate-return taps, which the current-capacity check counts as power copper.}

This is the one board that /fwe and /npie have worked on. /fwe wrote the first
stage of its firmware, which spins the motor with six-step commutation. It
builds with warnings treated as errors, its unit tests pass on the host, and
it runs on an emulated STM32G431 in Renode. /npie wrote a bring-up procedure
for the board and ran it once against a simulated bench, where it passed.
Neither result says anything about real hardware, because the board has not
been made.

\hsprovenance{Pictures from docs/showcase/renders, made as that folder's README says. Facts
from the board files of bldc-motor-driver in the boards repository.}
\newpage

% ----------------------------------------------------------------------
\subsection{The amplifier}

\board{stereo-class-d-amp}{stereo-class-d-amp}%
{2 layers \sep 53.6 $\times$ 36.8\,mm \sep 44 footprints}{verification, findings open}%
{A stereo class-D amplifier that runs from a single 12\,V supply and drives 4
to 8\,$\Omega$ speakers. Each output has a 10\,$\mu$H LC filter. The brief
asked hwde to reproduce a published reference design, so the two can be
compared.}

The brief named the TPA3116D2, but hwde used the TPA3118D2 instead. The 3116
only comes with its thermal pad on top, which needs a heatsink bolted on. The
3118 has the pad underneath and sheds its heat into the board, so hwde kept a
25 by 25\,mm ground area clear around it on both layers. That area is why the
board is bigger than the 44 by 34\,mm reference.

\board{stereo-class-d-amp-fcu}{stereo-class-d-amp, top copper}%
{F.Cu with net names}{}%
{The amplifier sits in the middle of a solid ground pour, with the left
channel's filter below it and the right channel's above. The input connector
is on the left edge and the speaker pads are on the right.}

\hsprovenance{Pictures from docs/showcase/renders, made as that folder's README says. Facts
from the board files of stereo-class-d-amp in the boards repository.}
\newpage

% ----------------------------------------------------------------------
\subsection{A card, a sensor node and a battery board}

These three were designed unattended on the same morning. The owner told hwde
to decide for itself, so hwde signed off its own stages and wrote down each
call it made, with the reason.

\setlength{\boardh}{6cm}
\board{nfc-card}{nfc-card}%
{2 layers \sep 88.9 $\times$ 50.8\,mm \sep 19 footprints}{package ready}%
{A business card you tap with a phone. The phone reads a web address from the
NFC tag and opens it, and while the phone is near, the tag takes power from
its field and lights a pattern of LEDs. There is no battery and no
microcontroller. The antenna is a coil printed around the edge of a 0.8\,mm
board, thin enough to carry in a wallet.}
\setlength{\boardh}{9cm}

\boardpair
{\board{esp32c3-node}{esp32c3-node}%
{2 layers \sep 26 $\times$ 38\,mm \sep 41 footprints}{package ready}%
{A small WiFi and Bluetooth sensor node built around an ESP32-C3 module. USB-C
powers it and also flashes it, through the chip's own USB, so there is no
USB-serial bridge. It charges a single LiPo cell and keeps running with or
without one, and a BME280 and a Qwiic connector share one
I\textsuperscript{2}C bus.}}
{\board{lipo-boost}{lipo-boost}%
{2 layers \sep 53.9 $\times$ 38.8\,mm \sep 33 footprints}{verification, findings open}%
{It charges a single LiPo cell from USB and makes 5\,V at 1\,A on a USB-A
socket from the cell. While USB is plugged in, the output keeps running from
USB as the cell charges. The brief asked for a board in the class of the
PowerBoost 1000C, designed from the datasheets rather than copied.}}

\hsprovenance{Pictures from docs/showcase/renders, made as that folder's README says. Facts
from the board files of nfc-card, esp32c3-node and lipo-boost in the boards
repository.}
\newpage

% ----------------------------------------------------------------------
\subsection{The small PD triggers}

\boardpair
{\board{pd-trigger-lite}{pd-trigger-lite}%
{2 layers \sep 25 $\times$ 15\,mm \sep 15 footprints}{package ready}%
{Plug in a USB-C charger and the board asks it for 9, 12, 15 or 20\,V, at up to
3\,A, and passes it to two output holes. One solder link picks the voltage.
hwde chose the CH224A over the CH224K the brief named, because its sense pin
takes the bus voltage directly. That removes a dropper resistor, a sense
resistor and the pull-ups.}}
{\board{pd-trigger-lite-dip}{pd-trigger-lite-dip}%
{2 layers \sep 25 $\times$ 21\,mm \sep 13 footprints}{package ready}%
{The same board with a five-way DIP switch instead of the solder links, which
adds 5\,V to the choices. You set exactly one switch before you plug it in.
The board grew from 15 to 21\,mm to make room for the switch.}}

Both are smaller versions of pd-trigger, the bench tool that was ordered and
made.

\hsprovenance{Pictures from docs/showcase/renders, made as that folder's README says. Facts
from the board files of pd-trigger-lite and pd-trigger-lite-dip in the boards
repository.}
\newpage

% ----------------------------------------------------------------------
\subsection{Ordered and fabricated}

lumina-carrier, the board at the start, is the carrier for a stage-lighting
fixture. It takes power over Ethernet through a PD controller and a 100\,V
buck, and it has an expansion connector that a fixture-specific daughter board
plugs into. The connector's pinout is frozen, because another board depends on
it. The other board that was ordered is a bench tool.

\board{pd-trigger}{pd-trigger}%
{2 layers \sep 48 $\times$ 30\,mm \sep 30 footprints}{ordered and made}%
{A USB-C Power Delivery sink asks a charger for 5, 9, 12, 15 or 20\,V, and you
pick which on the board. The rail comes out on a screw terminal rated for
5\,A. The board converts nothing and just passes the rail through.}

\hsprovenance{Pictures from docs/showcase/renders, made as that folder's README says. Facts
from the board files of pd-trigger in the boards repository.}
\newpage

% ----------------------------------------------------------------------
\subsection{Power and RF}

\boardpair
{\board{sbuck-5v3a}{sbuck-5v3a}%
{4 layers \sep 50 $\times$ 40\,mm \sep 44 footprints}{package ready}%
{A synchronous buck that takes 7--18\,V in and holds 5.0\,V $\pm$2\% at 3\,A
out. It is an open-frame power module.}}
{\board{rf-term-150w}{rf-term-150w}%
{2 layers \sep 26 $\times$ 20\,mm \sep 6 footprints}{package ready}%
{A 50\,$\Omega$, 150\,W dummy load head for DC to 25\,MHz. The element bolts to
a heatsink you supply, so the board is only the RF launch and the mounting. It
looks almost bare because none of its six footprints has a 3D model, so the
connector and the element are not drawn.}}

\newpage

\board{rf-de-20m}{rf-de-20m}%
{4 layers \sep 120 $\times$ 80\,mm \sep 80 footprints}{routing, findings open}%
{A 20\,MHz Class E RF power stage that puts 200\,W into 50\,$\Omega$ from a
40\,V bus. A GaN gate driver buffers the PWM drive from an SMA input, and one
eGaN FET switches. A resonant tank and an L-match turn the 4.6\,$\Omega$ load
line into 50\,$\Omega$ at the output SMA. The owner fixed the topology and the
two core parts after a costed trade study, and hwde did the rest.}

\hsprovenance{Pictures from docs/showcase/renders, made as that folder's README says. Facts
from the board files of sbuck-5v3a, rf-term-150w and rf-de-20m in the boards
repository.}
\newpage

% ----------------------------------------------------------------------
\subsection{Lighting}

\setlength{\boardh}{6cm}
\board{lumina-par}{lumina-par}%
{4 layers \sep 100 $\times$ 80\,mm \sep 158 footprints}{manufacturability, findings open}%
{The RGBW PAR daughter board that stacks on lumina-carrier. It takes the
carrier's frozen connector as given. Its power budget depends on the
power-over-Ethernet class the fixture is granted, so hwde worked it out again
here. The silkscreen warns that the board floats at the power-over-Ethernet
potential, because an earthed probe breaks the fixture's power negotiation.}
\setlength{\boardh}{9cm}

\boardpair
{\board{lumina-par-flat}{lumina-par, straight down}%
{top-down view}{}%
{The same board from straight above, so the routing shows.}}
{\board{lumina-carrier-bottom}{lumina-carrier, underneath}%
{bottom view}{}%
{Every part on the carrier is on the top side, so the back carries routing and
vias and nothing else.}}

\hsprovenance{Pictures from docs/showcase/renders, made as that folder's README says. Facts
from the board files of lumina-par in the boards repository.}
\newpage

% ----------------------------------------------------------------------
\subsection{Microcontroller boards}

\setlength{\boardh}{6cm}
\board{g0-sense}{g0-sense}%
{2 layers \sep 35.8 $\times$ 28.3\,mm \sep 30 footprints}{package ready}%
{A temperature and humidity node that runs off USB-C. An LDO makes 3.3\,V from
the 5\,V supply, and an STM32G030 reads a Sensirion SHT4x over
I\textsuperscript{2}C. The same bus comes out on a Qwiic connector, the
readings go out on a UART header, and you program it over SWD. hwde designed
this one from end to end in a single unattended run.}
\setlength{\boardh}{9cm}

\boardpair
{\board{usb-buck}{usb-buck}%
{4 layers \sep 50 $\times$ 40\,mm \sep 28 footprints}{package ready}%
{An STM32F103 board that talks USB at full speed. The Micro-B socket carries
both power and data, and there is an LED, a button and an SWD header.}}
{\board{stm32-blinky}{stm32-blinky}%
{2 layers \sep 50 $\times$ 40\,mm \sep 20 footprints}{package ready}%
{The simplest useful board here: an STM32F103, an 8\,MHz crystal, one LED, an
LDO and an SWD header.}}

\hsprovenance{Pictures from docs/showcase/renders, made as that folder's README says. Facts
from the board files of g0-sense, usb-buck and stm32-blinky in the boards
repository.}
\newpage

% ----------------------------------------------------------------------
\subsection{Building blocks}

These five boards each carry a single block and nothing else. They are meant
to be studied and measured on a bench, so they leave out protection,
filtering, indicators and spare rails on purpose.

\boardpair
{\board{bb-buck}{bb-buck}%
{2 layers \sep 35 $\times$ 25\,mm \sep 20 footprints}{package ready}%
{18--30\,V in and 5\,V at 2\,A out, on a screw terminal. It carries the
converter and exactly what its datasheet requires.}}
{\board{bb-ldo}{bb-ldo}%
{2 layers \sep 34.7 $\times$ 34.7\,mm \sep 5 footprints}{package ready}%
{5\,V in and 3.3\,V at 500\,mA out. It stays in regulation in still air on the
board's own copper.}}

\boardpair
{\board{bb-mcu}{bb-mcu}%
{2 layers \sep 34.8 $\times$ 22.3\,mm \sep 14 footprints}{package ready}%
{One MCU with a 3.3\,V rail coming in, its debug header and four GPIO. There is
no regulator and no second rail.}}
{\board{bb-adc}{bb-adc}%
{2 layers \sep 54.8 $\times$ 34.9\,mm \sep 25 footprints}{package ready}%
{0--5\,V in on a screw terminal and 12 bits at up to 10\,kSa/s out on a header.
It gives up speed for DC accuracy.}}

\newpage

\boardpair
{\board{bb-amp}{bb-amp}%
{2 layers \sep 48 $\times$ 28.3\,mm \sep 14 footprints}{package ready}%
{A 0--20\,mV bridge signal in and 0--3.3\,V out, from DC to 1\,kHz.}}
{}

\hsprovenance{Pictures from docs/showcase/renders, made as that folder's README says. Facts
from the board files of bb-buck, bb-ldo, bb-mcu, bb-adc and bb-amp in the
boards repository.}

Two more boards stopped before layout and have no PCB to show: lumina-strobe
(LUMINA strobe daughter board) and buck-5v3a (12\,V to 5\,V / 3\,A synchronous
buck). They are in the boards repository at the stage they reached.

\newpage
